\noindent{\bf Retrieval-Augmented Generation.}
{
Retrieval-Augmented Generation (RAG) augments the input space of LMs with retrieved text passages~\citep{guu2020retrieval,lewis2020retrieval}, leading to large improvements in knowledge-intensive tasks after fine-tuning or used with off-the-shelf LMs~\citep{ram2023context}. 
A more recent work~\citep{luo2023sail} instruction-tunes an LM with a fixed number of retrieved passages prepended to input, or pre-train a retriever and LM jointly, followed by few-shot fine-tuning on task datasets~\citep{izacard2022few}. }
While prior work often retrieves only once at the beginning, \cite{jiang2023active} propose to
adaptively retrieve passages for generation on top of a proprietary LLM or ~\cite{schick2023toolformer} train an LM to generate API calls for named entities. 
Yet, the improved task performance of such approaches often comes at the expense of runtime efficiency~\citep{mallen2022not}, robustness to irrelevant context~\citep{pmlr-v202-shi23a}, and lack of attributions~\citep{liu2023evaluating,gao2023enabling}. 
We introduce a method to train an arbitrary LM to learn to use retrieval \textit{on-demand} for diverse instruction-following queries and introduce controlled generation guided by reflections tokens to further improve generation quality and attributions. 

\noindent{\bf Concurrent RAG work. }
{
A few concurrent works\footnote{All work is arXived within a week of this preprint.} on RAG propose new training or prompting strategies to improve widely-adopted RAG approaches. 
\cite{lin2023radit} fine-tune both the retriever and LM on instruction-tuning datasets in two steps.
While we also train our model on diverse instruction-following datasets, \model enables retrieval on demand and selection of the best possible model output via fine-grained self-reflection, making it widely applicable and more robust and controllable. 
\citet{Yoran2023MakingRL} use a natural language inference model and \citet{xu2023recomp} use a summarization model to filter out or compress retrieved passages before using them to prompt the LM to generate the output.
\model processes passages in parallel and filters out irrelevant ones through self-reflection, without relying on external models at inference. Moreover, our self-reflection mechanism also evaluates other aspects of the model output quality including factuality.
LATS~\citep{zhou2023language} prompt off-the-shelf LMs to search for relevant information for question answering tasks and to generate with tree search, guided by LM-generated value scores. 
While their value function simply indicates an overall score of each generation, \model trains to an arbitrary LM to learn to generate fine-grained self-reflection and customizable inference. 
% to guide the generation.
%While LATS uses web search for QA tasks, they do not introduce fine-grained multi-aspect feedback to guide search as in \model. 
}

% feature by carefully prompting a proprietary LM for QA and summarization tasks~\citep{jiang2023active} or by augmenting pre-training data with API calls~\citep{schick2023toolformer}. 
% This work introduces an end-to-end training of an adaptive retrieval-augmented LM that learns to flexibly retrieve when necessary while generation.
% We train a general-purpose retrieval-augmented instruction-following LM that can be widely applicable and customized for diverse downstream tasks, and we further propose a new inference-time algorithm built upon our training method to ensure that model generation is supported by the retrieved text.  


\noindent{\bf Training and generating with critics.}
Training LLMs with reinforcement learning (e.g., Proximal Policy Optimization or PPO; \citealt{schulman2017proximal}) from human feedback (RLHF) has proven effective in aligning LLMs with human preferences~\citep{ouyang2022training}. 
% \cite{menick2022teaching} trains retrieval-based QA systems using finer-grained feedback, including feedback on retrieved evidence as well as attributions, instead of overall preferences only as in those prior studies.
\citet{wu2023fine} introduce fine-grained RLHF with multiple reward models. 
% and demonstrates the effectiveness of giving feedback from multiple aspects. 
Though our work also studies fine-grained critique on retrieval and generation, we train our target LM on task examples augmented with reflection tokens from a critic 
% \el{I think we need to be consistent when to use critic/critique?} 
model offline, with a far lower training cost compared to RLHF. In addition, reflection tokens in \model enable controllable generation at inference, while RLHF focuses on human preference alignment during training.
Other works use general control tokens to guide LM generation~\citep{lu2022quark,korbak2023pretraining}, while \model uses reflection tokens to decide the need for retrieval and to self-evaluate generation quality.
%Embedding special tokens into an LM training corpus and training LMs using standard supervised training objectives to control generation at test time has been recently studied~\citep{lu2022quark,korbak2023pre-training}. 
% , removing the necessity of relying on such external reward modules at training or test time. 
% by training the target LM on an augmented corpus created by our critique model offline. 
% Embedding special tokens into an LM training corpus and training LMs using standard objective has been recently studied~\citep{lu2022quark,korbak2023pre-training}, by conditioning LM outputs on precedding special tokens. 
\cite{xie2023decomposition} propose a self-evaluation-guided decoding framework, but they focus only on reasoning tasks with one evaluation dimension (reasoning path consistency) and without retrieval.
%Prior work explores prompting techniques on the top of proprietary LLMs to enable LMs to self-evaluate intermediate reasoning steps~\citep{xie2023decomposition}. 
%Yet, it has not yet explored what type of general critique feedback should be considered for general tasks, and moreover, whether a smaller LM can reliably conduct self-evaluation or retrieval. 
% 
Recent work on LLM refinement \citep{dhuliawala2023chainofverification, madaan2023selfrefine, paul2023refiner} prompts a model to generate task output, natural language feedback and refined task output iteratively, but at the cost of inference efficiency.